\subsection*{Training corpus}

The dated corpus comprises 25 units spanning 760--167~BCE and 146{,}059
tokens: 23 training units and two units held out of both models
(Table~\ref{tab:corpus}).  Date anchoring uses explicit synchronisms with
datable ancient Near Eastern events wherever these are available.  Amos~1:1
places that book in the reigns of Uzziah and Jeroboam~II ($\sim$760~BCE);
Zechariah~1--8 and Haggai carry explicit Darius regnal dates (518 and
520~BCE); Ezekiel supplies thirteen regnal dates checkable against
Babylonian records; Daniel's Hebrew chapters~(1,~8--12) are anchored to
167~BCE by the Maccabean crisis.  Daniel is restricted to those Hebrew
chapters; the Aramaic chapters~(2--7) are a different language and are
excluded.

Not every unit is anchored this securely, and the corpus does not pretend
otherwise.  Where a date rests on literary-critical rather than
synchronistic grounds, $\sigma$ is widened accordingly: Joel and
Zechariah~9--14 carry $\sigma = 150$~yr, Ecclesiastes 120~yr, Isaiah~56--66
100~yr, and Obadiah 80~yr, against 10--30~yr for the synchronistically
anchored texts.  This is a substantial widening relative to earlier versions
of this corpus, and it is deliberate: a narrow prior on a literary-critical
date imports the very assumption the analysis is meant to test.  The prior
sensitivity analysis below quantifies what that assumption is worth.

Two units are held out of both models and excluded from all feature
screening: the Jeremiah oracular stratum (chs.~1--6, 8--10, 12--16, 19--20,
22--23, 30--31, 46--51) and Haggai.  Whole-book Jeremiah is excluded from
training entirely, because it contains the oracular stratum and would
otherwise leak the holdout into the training set.  The Deuteronomistic prose
stratum of Jeremiah (chs.~7, 11, 17--18, 21, 24--29, 32--45, 52) is treated
as an undated target and never trained on.  Habakkuk and Daniel are retained
in MLE-MVN training but held out of HB-VI.

A note on what ``held out'' must mean.  In the analyses reported below, a
held-out unit contributes nothing to feature screening, to standardisation,
to regression fitting, or to its own prior.  The last of these is not a
formality.  A model that receives a text's scholarly date as a prior mean
will return approximately that date whatever the text contains, and will do
so most emphatically for the texts whose dates are best established --- that
is, precisely the texts on which a validation exercise appears most
convincing.  Section~\nameref{sec:validation} quantifies this effect for the
present corpus.  All out-of-sample results in this paper use an agnostic
prior.

S5~Table lists word counts for all units, including the Documentary
Hypothesis source composites, which are excluded from both training and
validation and are dated solely from the model.
